\capitulo{2}{Objetivos del proyecto}

El objetivo principal de este proyecto es desarrollar una aplicación web docente que facilite el aprendizaje del algoritmo C4.5 y de los conceptos necesarios para comprender la construcción de árboles de decisión. La herramienta está orientada a usuarios sin conocimientos previos, por lo que debe presentar los contenidos de forma clara, visual e interactiva.

La aplicación debe permitir al usuario trabajar con conceptos básicos como la entropía y la entropía condicional, introduciendo valores personalizados y visualizando los resultados obtenidos. Asimismo, debe ofrecer una sección dedicada a la elección de nodos, en la que se expliquen y calculen los elementos utilizados por el algoritmo C4.5, como los umbrales, la ganancia de información, el \textit{split info} y el \textit{gain ratio}.

Otro objetivo fundamental es proporcionar una simulación paso a paso de la construcción de un árbol de decisión mediante el algoritmo C4.5. Para ello, el usuario debe poder utilizar conjuntos de datos de ejemplo o cargar archivos CSV propios, siempre que cumplan las condiciones de validación establecidas por la aplicación. Durante la simulación, se deben mostrar de forma sincronizada el árbol generado, el conjunto de datos utilizado y los cálculos realizados en cada iteración.

Además, la herramienta debe incluir una simulación del proceso de poda del árbol, permitiendo observar cómo se evalúan los subárboles, qué errores se calculan y en qué casos se decide sustituir un subárbol por una hoja. De esta forma, el usuario puede comprender tanto la generación inicial del árbol como su posterior simplificación.

Por último, la aplicación debe ofrecer una interfaz sencilla, intuitiva y adaptable a distintos tamaños de pantalla. La representación visual del árbol debe mantenerse legible incluso con conjuntos de datos relativamente grandes, incorporando herramientas de ampliación, lupa, paneles retráctiles, resaltado del nodo actual y ayudas contextuales que faciliten el seguimiento de la simulación.

